% ============================================================
%  TEIA Cognitive Router — arXiv Submission
%  "Deterministic and Auditable Routing for Artificial
%   Intelligence in Regulated Environments"
%
%  Compile: pdflatex teia_arxiv_paper.tex  (run twice)
%  arXiv categories: cs.CR, cs.AI, cs.LG
% ============================================================

\documentclass[11pt,a4paper]{article}

% ── Packages ─────────────────────────────────────────────────────────────────
\usepackage[margin=1in]{geometry}
\usepackage{amsmath,amssymb,amsthm}
\usepackage{booktabs}
\usepackage{hyperref}
\usepackage{xcolor}
\usepackage{listings}
\usepackage{microtype}
\usepackage{parskip}
\usepackage{enumitem}
\usepackage{tabularx}
\usepackage{array}

% ── Theorem environments ──────────────────────────────────────────────────────
\newtheorem{theorem}{Theorem}[section]
\newtheorem{corollary}[theorem]{Corollary}
\newtheorem{lemma}[theorem]{Lemma}
\newtheorem{definition}[theorem]{Definition}

% ── Code listing style ────────────────────────────────────────────────────────
\lstset{
  basicstyle=\ttfamily\small,
  breaklines=true,
  frame=single,
  backgroundcolor=\color{gray!8},
  keywordstyle=\color{blue!70!black}\bfseries,
  commentstyle=\color{green!50!black},
  stringstyle=\color{orange!70!black},
  showstringspaces=false,
  columns=flexible,
}

% ── Hyperref setup ────────────────────────────────────────────────────────────
\hypersetup{
  colorlinks=true,
  linkcolor=blue!60!black,
  citecolor=green!50!black,
  urlcolor=blue!70!black,
  pdftitle={Deterministic and Auditable Routing for Artificial Intelligence in Regulated Environments},
  pdfauthor={Felippe Barcelos},
  pdfkeywords={LLM routing, compliance, determinism, cryptographic audit, EU AI Act},
}

% ── Math shortcuts ────────────────────────────────────────────────────────────
\newcommand{\Et}{E(t)}
\newcommand{\sha}{\texttt{SHA-256}}
\newcommand{\Ai}{A_i}
\newcommand{\Bi}{B_i}

% ─────────────────────────────────────────────────────────────────────────────
\begin{document}

\title{\textbf{Deterministic and Auditable Routing for Artificial\\
Intelligence in Regulated Environments}}

\author{
  Felippe Barcelos\\
  TEIA Project (Independent)\\
  \texttt{felippe.barcelos10@gmail.com}\\[4pt]
  \small\url{https://github.com/felippebarcelos/teia-omega-awakening}\\
  \small\texttt{pip install teia-cognitive-router}
}

\date{June 2026 \quad|\quad Protocol P51.0}

\maketitle

% ── Abstract ─────────────────────────────────────────────────────────────────
\begin{abstract}
Large Language Model (LLM) routing systems---software that selects whether a
given prompt should be processed by a local small model, a quantized hybrid
model, or a frontier cloud LLM---are increasingly deployed in regulated
industries such as financial services, healthcare, and government contracting.
Existing machine-learning-based routers produce non-reproducible decisions: the
same prompt submitted on different dates may receive different routing verdicts
due to weight drift from model retraining, dependency on third-party embedding
services, or non-deterministic sampling. This property is incompatible with the
audit requirements of the EU AI Act (High-Risk System provisions, Articles
12--14 and Annex~IV), GDPR Article~22, SEC/FINRA algorithmic accountability
rules, and the HIPAA Security Rule audit controls standard.

We present \textbf{TEIA Cognitive Router}, a deterministic, compliance-first
LLM routing system that produces a cryptographically sealed, reproducible
routing decision for every prompt. The routing formula is a fixed linear
combination of six measurable text features---token count, vocabulary diversity,
reasoning verb density, data operation density, structural complexity, and
constraint density---calibrated at Protocol P40.0 and version-frozen thereafter.
Every decision is accompanied by a \sha\ body seal, a human-readable rationale
string, and a Merkle-chained temporal anchor that links each decision to its
predecessor in an append-only audit ledger that can be externally notarized via
RFC~3161 Trusted Timestamping.

We evaluate TEIA on MT-Bench~\cite{zheng2023mtbench} (80 questions, 8
categories) and report \textbf{99.6\% quality retention} at \textbf{16.3\% cost
reduction} versus routing all prompts to a frontier LLM. The system achieves
P99 routing latency below 1\,ms and throughput exceeding 15{,}000 decisions per
second on commodity hardware, with zero network calls and a memory footprint
below 5\,MB. We prove formally that the routing formula satisfies the
\emph{Write\,==\,Read invariant}---identical input always produces an identical
\sha\ seal---and describe the cryptographic properties of the temporal chain
that make retroactive modification mathematically detectable.

\textbf{Keywords:} LLM routing, compliance, determinism, cryptographic audit,
EU AI Act, GDPR, RFC 3161, AI governance.
\end{abstract}

\tableofcontents
\newpage

% ─────────────────────────────────────────────────────────────────────────────
\section{Introduction}
\label{sec:intro}

\subsection{Motivation}

The widespread deployment of large language models in regulated industries has
created a cost--quality optimization problem: frontier models such as GPT-4 and
Claude~3 Opus deliver high-quality outputs but cost orders of magnitude more per
token than smaller local models. A routing layer that selects the cheapest model
capable of satisfying a given request can reduce infrastructure costs
substantially.

Regulated industries, however, face an additional constraint that purely
cost-focused routing systems do not address: \emph{every AI-assisted decision
affecting a user or business outcome must be explainable, reproducible, and
auditable.} A compliance officer who asks ``why did this prompt route to GPT-4
on March~15 instead of the local model?'' must receive a mathematically precise
and independently verifiable answer. These requirements are not operational
preferences---they are legal mandates codified in multiple overlapping
regulatory frameworks.

\subsection{The Compliance Failure of ML-Based Routing}
\label{sec:compliance-failure}

Current state-of-the-art LLM routers employ machine learning as their decision
mechanism. Representative systems include FrugalGPT~\cite{chen2023frugal},
which uses a cascading strategy with a learned escalation policy; Hybrid
LLM~\cite{ding2024hybrid}, which trains a binary routing classifier; and
RouteLLM~\cite{ong2024routellm}, which frames routing as a preference learning
problem. These systems achieve strong cost--quality Pareto improvements but
share a structural property that disqualifies them from regulated deployments:
their routing decisions are not reproducible across model versions.

\begin{table}[h]
\centering
\caption{Compliance failures of ML-based LLM routers.}
\label{tab:compliance-failures}
\begin{tabularx}{\textwidth}{lX}
\toprule
\textbf{ML Router Property} & \textbf{Regulatory Consequence} \\
\midrule
Weights updated at retraining & Same prompt routes differently next quarter
  --- historical audit trail is broken \\[4pt]
Decision depends on embedding distance & Identical text produces different
  embeddings across model versions \\[4pt]
Post-hoc explanation only & Cannot satisfy ``right to explanation'' at decision
  time \\[4pt]
Routing service is a network dependency & Reproducibility depends on
  third-party uptime and API versioning \\[4pt]
GPU/model inference in routing path & Adds 50--500\,ms latency; routing failure
  cascades to service outage \\
\bottomrule
\end{tabularx}
\end{table}

The fundamental problem is that \emph{accuracy without reproducibility is
insufficient for compliance.} A routing decision that cannot be independently
re-derived from the original input, using only the original routing logic,
cannot satisfy the legal burden of proof required by EU AI Act Annex~IV,
GDPR Article~22, or SEC/FINRA record reconstruction requirements.

\subsection{Contributions}
\label{sec:contributions}

This paper makes the following contributions:

\begin{enumerate}[leftmargin=2em]
  \item \textbf{A deterministic 6-axis heuristic routing formula} (Section~\ref{sec:formula})
    that maps any text string to a routing tier in $O(|t|)$ time with zero ML
    dependencies, zero network calls, and sub-millisecond latency.

  \item \textbf{A cryptographic audit seal} (Section~\ref{sec:seal}) providing a
    \sha\ commitment to every routing decision body. The Write\,==\,Read
    invariant guarantees identical input produces an identical seal, enabling
    independent verification.

  \item \textbf{A Merkle-chained temporal ledger} (Section~\ref{sec:chain})
    that links routing decisions into an append-only sequence. Modification,
    deletion, or reordering of any entry is mathematically detectable.

  \item \textbf{An RFC~3161 external notarization module} (Section~\ref{sec:rfc3161})
    that closes the retroactive forgery window by anchoring the chain to an
    externally-witnessed, cryptographically-signed timestamp.

  \item \textbf{An empirical evaluation on MT-Bench} (Section~\ref{sec:eval})
    demonstrating 99.6\% quality retention at 16.3\% cost reduction with
    P99 routing latency below 1\,ms.

  \item \textbf{A regulatory alignment analysis} (Section~\ref{sec:regulatory})
    mapping each system property to specific EU AI Act, GDPR, HIPAA, and
    SEC/FINRA requirements.
\end{enumerate}

% ─────────────────────────────────────────────────────────────────────────────
\section{Background and Related Work}
\label{sec:background}

\subsection{LLM Cost--Quality Routing}

LLM routing research aims to characterize and exploit the cost--quality
trade-off in large language model deployments.

\textbf{FrugalGPT}~\cite{chen2023frugal} introduced a cascading framework in
which a prompt is submitted to progressively larger and more expensive models
until a quality threshold is satisfied. The escalation policy is learned from
historical data and achieves cost savings up to 98\% with minimal quality
degradation. However, the learned policy is not fixed and changes whenever the
quality estimator is retrained.

\textbf{Hybrid LLM}~\cite{ding2024hybrid} trains a small binary routing
classifier on human preference data to predict whether a small LLM will handle
a given prompt adequately. It achieves strong Pareto improvements but inherits
the non-reproducibility of any ML-based classifier.

\textbf{RouteLLM}~\cite{ong2024routellm} frames routing as a preference
learning problem, training the router to maximize preference satisfaction under
a cost budget. RouteLLM introduces evaluation protocols for measuring routing
efficiency and achieves favorable Pareto curves on standard benchmarks. Like
prior ML-based routers, its decisions are not reproducible across model
checkpoints.

\textbf{LLM-Blender}~\cite{jiang2023blender} proposes an ensemble approach that
ranks and fuses outputs from multiple LLMs. While primarily an output-selection
strategy, it shares the ML-based non-reproducibility property.

TEIA differs from all four systems by optimizing for compliance reproducibility
as a primary objective, with cost efficiency as a secondary objective.

\subsection{Heuristic and Rule-Based Routing}

Rule-based routing systems predate ML-based approaches. Simple heuristics ---
route short prompts to local models, route long prompts to cloud models --- have
obvious compliance advantages but perform poorly on the cost--quality Pareto
curve because prompt length is a weak predictor of required reasoning
complexity.

TEIA's 6-axis formula is a structured heuristic: more expressive than a single
rule yet fully deterministic and auditable. The formula's weights are fixed at
Protocol P40.0, giving it the compliance properties of a rule-based system
while approaching the expressiveness of a shallow classifier.

\subsection{Cryptographic Audit Logging}

Certificate Transparency~\cite{laurie2013ct} (RFC~6962) maintains a publicly
auditable Merkle tree of all issued X.509 certificates. Each log entry commits
to the hash of the previous entry; presenting different views to different
observers is detectable by cross-checking inclusion proofs. TEIA's temporal
chain follows the same structural pattern.

\textbf{RFC~3161 Trusted Timestamping}~\cite{adams2001rfc3161} extends hash
commitments with external time witnesses. A Time Stamp Authority (TSA) receives
a hash, signs it with the TSA's private key, and returns a signed
TimeStampResponse (TSR). The TSR is legally equivalent to a notarized timestamp
under eIDAS Regulation (EU) No~910/2014, Article~41.

\subsection{AI Governance and Auditability}

Doshi-Velez and Kim~\cite{doshi2017interpretable} distinguish
\emph{application-grounded} evaluation (testing in real use cases) from
\emph{functionally-grounded} evaluation (using a proxy). TEIA's
\texttt{routing\_rationale} field targets functional grounding: the explanation
is produced by the same arithmetic used to make the decision, not by a
post-hoc model.

The EU AI Act~\cite{euaiact2024} represents the most comprehensive regulatory
codification of AI auditability requirements. High-Risk AI systems must maintain
technical documentation (Annex~IV) including descriptions of the measures
ensuring accuracy, robustness, and cybersecurity.

% ─────────────────────────────────────────────────────────────────────────────
\section{TEIA Cognitive Router: System Design}
\label{sec:formula}

\subsection{Design Constraints}

The routing formula is designed under three constraints:

\begin{description}[leftmargin=2em]
  \item[C1 --- Closed-form determinism.] The routing decision is a mathematical
    function of the input text string. No learned weights, external service
    calls, random sampling, or system clock access occur during routing.

  \item[C2 --- Version stability.] Formula parameters are pinned to Protocol
    P40.0. Any modification constitutes a new protocol version with a new
    identifier, embedded in every routing decision.

  \item[C3 --- Human interpretability at decision time.] Every feature is a
    directly measurable, interpretable property of the text. The
    \texttt{routing\_rationale} field states in natural language which features
    drove the verdict.
\end{description}

\subsection{Feature Definitions}
\label{sec:features}

Let $t$ denote an input text string with character count $|t|$.
Let $W(t)$ be the token list obtained by splitting $t$ on non-alphanumeric
boundaries (discarding tokens of length $\leq 1$, lowercased).
Let $n_w = \max(1, |W(t)|)$ and $n_u = |\text{unique}(W(t))|$.

We define six normalized features $f_i : \Sigma^* \to [0,1]$.

\begin{definition}[Token Score, $f_1$]
\begin{equation}
f_1(t) = \min\!\left(1.0,\; \frac{|t|}{16000}\right)
\label{eq:f1}
\end{equation}
The reference length $16{,}000$ characters corresponds to approximately 4{,}000
tokens at four characters per token.
\end{definition}

\begin{definition}[Vocabulary Diversity, $f_2$]
Let $r = n_u / n_w$ be the unique-token ratio.
\begin{equation}
f_2(t) = \max\!\left(0.0,\; \min\!\left(1.0,\; \frac{r - 0.30}{0.60}\right)\right)
\label{eq:f2}
\end{equation}
\end{definition}

\begin{definition}[Reasoning Verb Density, $f_3$]
Let $V_R$ be a fixed vocabulary of 44 analytical verbs in English and Portuguese
(analyze, synthesize, evaluate, compare, hypothesize, \ldots). Let
$c_3 = |\{v \in V_R : v \in \text{lowercase}(t)\}|$.
\begin{equation}
f_3(t) = \min\!\left(1.0,\; \frac{c_3 / n_w}{0.04}\right)
\label{eq:f3}
\end{equation}
\end{definition}

\begin{definition}[Data Operation Score, $f_4$ --- inverted]
Let $V_D$ be a fixed vocabulary of 45 data manipulation verbs (list, extract,
format, sort, filter, \ldots). Let $c_4 = |\{v \in V_D : v \in \text{lowercase}(t)\}|$.
\begin{equation}
f_4(t) = \max\!\left(0.0,\; 1.0 - \min\!\left(1.0,\; \frac{c_4 / n_w}{0.04}\right)\right)
\label{eq:f4}
\end{equation}
High data verb density \emph{reduces} entropy, signaling a task within local
model capability.
\end{definition}

\begin{definition}[Structural Complexity, $f_5$]
Let $c_b$ = code block count, $b_p$ = bullet point count,
$h_d$ = header count, $l_c$ = line count.
\begin{equation}
f_5(t) = \min\!\left(1.0,\; \frac{3\,c_b + 0.5\,b_p + h_d + l_c/20}{15.0}\right)
\label{eq:f5}
\end{equation}
\end{definition}

\begin{definition}[Constraint Density, $f_6$]
Let $P_C$ be a set of 26 logical connectives (but, however, unless, given~that,
\ldots). Let $c_6 = \sum_{p \in P_C}\text{count}(p, \text{lowercase}(t))$ and
$q_6 = \text{count}(\text{`?`}, t)$.
\begin{equation}
f_6(t) = \min\!\left(1.0,\; \frac{(c_6 + 0.5\,q_6) / n_w}{0.10}\right)
\label{eq:f6}
\end{equation}
\end{definition}

\subsection{Composite Semantic Entropy Score}
\label{sec:entropy}

\begin{definition}[Semantic Entropy Score]
\begin{equation}
E(t) = \sum_{i=1}^{6} w_i \cdot f_i(t)
     = 0.20\,f_1 + 0.15\,f_2 + 0.30\,f_3 + 0.15\,f_4 + 0.10\,f_5 + 0.10\,f_6
\label{eq:entropy}
\end{equation}
\end{definition}

Since $\sum_{i=1}^{6} w_i = 1.00$ and $f_i \in [0,1]$ for all $i$, we have
$E(t) \in [0,1]$.

\begin{table}[h]
\centering
\caption{Feature weights and interpretive roles (Protocol P40.0).}
\label{tab:weights}
\begin{tabular}{llrc}
\toprule
\textbf{Feature} & \textbf{Role} & \textbf{Weight} & \textbf{Direction} \\
\midrule
$f_1$ Token Score & Information volume proxy & 0.20 & $\uparrow$ \\
$f_2$ Vocabulary Diversity & Conceptual breadth proxy & 0.15 & $\uparrow$ \\
$f_3$ Reasoning Verb Density & \textbf{Dominant signal} --- analytical complexity & 0.30 & $\uparrow$ \\
$f_4$ Data Operation Score & Simple task signal (inverted) & 0.15 & $\downarrow$ \\
$f_5$ Structural Complexity & Multi-part query approximation & 0.10 & $\uparrow$ \\
$f_6$ Constraint Density & Multi-condition requirement signal & 0.10 & $\uparrow$ \\
\bottomrule
\end{tabular}
\end{table}

\subsection{Routing Decision}
\label{sec:routing-decision}

TEIA supports two routing policies:

\medskip
\noindent\textbf{Compliance-safe mode} (default, $\tau_{cs} = 0.20$):
\begin{equation}
\text{routing}_{cs}(t) =
\begin{cases}
\text{Local} & \text{if } E(t) < \tau_{cs} \\
\text{Cloud} & \text{otherwise}
\end{cases}
\label{eq:routing-cs}
\end{equation}

\noindent\textbf{Max-savings mode} (opt-in, $\tau_L = 0.35$, $\tau_H = 0.65$):
\begin{equation}
\text{routing}_{ms}(t) =
\begin{cases}
\text{Local}  & \text{if } E(t) < \tau_L \\
\text{Hybrid} & \text{if } \tau_L \leq E(t) < \tau_H \\
\text{Cloud}  & \text{if } E(t) \geq \tau_H
\end{cases}
\label{eq:routing-ms}
\end{equation}

The routing confidence signal is:
\begin{equation}
\text{confidence}(E) =
\begin{cases}
\text{high}   & \text{if } \min(|E - \tau_L|, |E - \tau_H|) > 0.12 \\
\text{medium} & \text{if } \min(|E - \tau_L|, |E - \tau_H|) > 0.05 \\
\text{low}    & \text{otherwise}
\end{cases}
\end{equation}

\subsection{Formal Determinism Proof}
\label{sec:determinism}

\begin{theorem}[Write\,==\,Read Invariant]
\label{thm:wr}
For any text string $t \in \Sigma^*$, the function $\mathtt{route}(t)$ returns
identical output on every invocation with identical input. Consequently,
$\sha(\mathtt{canonical\_json}(\mathtt{route}(t)))$ is identical for every
invocation.
\end{theorem}

\begin{proof}
Inspect the computation graph of $\mathtt{route}(t)$:

\begin{enumerate}
  \item Features $f_1$--$f_6$ are computed by closed-form arithmetic over the
    character sequence of $t$: integer/float arithmetic, character counting,
    set membership tests against compile-time-constant vocabularies, and regular
    expression matching with fixed patterns. None of these operations involves
    random number generation, I/O, system calls, or external network access.

  \item $E(t)$ (Eq.~\ref{eq:entropy}) is a fixed linear combination of
    $f_1$--$f_6$ with compile-time-constant weights. It is therefore a
    deterministic function of $t$.

  \item The routing decision (Eqs.~\ref{eq:routing-cs},\ref{eq:routing-ms}) is a
    piecewise-constant function of $E(t)$ with fixed thresholds. It is therefore
    a deterministic function of $t$.

  \item All remaining fields of the routing result (rationale, confidence, GPU
    economics) are deterministic functions of the routing decision and feature
    scores.

  \item The canonical JSON serializer uses \texttt{sort\_keys=True},
    \texttt{separators=(',',':')}, and \texttt{ensure\_ascii=False} ---
    eliminating key-order non-determinism and formatting variation.

  \item \sha\ is a deterministic hash function.
\end{enumerate}

Therefore $\mathtt{route}(t)$ is a pure function, and the composition
$\sha(\mathtt{canonical\_json}(\mathtt{route}(t)))$ is likewise a pure function
of $t$. \qed
\end{proof}

\begin{corollary}
\label{cor:auditor}
Any compliance auditor in possession of the original prompt text $t$ and the
TEIA Cognitive Router source code at Protocol~P40.0 can independently
re-derive the routing decision and verify its \sha\ seal without any network
access or special infrastructure.
\end{corollary}

% ─────────────────────────────────────────────────────────────────────────────
\section{Cryptographic Audit Architecture}

\subsection{Routing Decision Seal}
\label{sec:seal}

After computing $\mathtt{route}(t)$, the $\mathtt{seal}()$ operation:
\begin{enumerate}
  \item Serializes the routing result to canonical JSON (sorted keys, no
    whitespace, UTF-8 encoding).
  \item Computes $B = \sha(\text{canonical\_json\_bytes})$ --- the body hash.
  \item Attaches an \texttt{audit\_seal} block: $\{\texttt{sha256}: B,\
    \texttt{algorithm}: \text{``SHA-256''}, \texttt{scope}:
    \text{``canonical\_json\_utf8''}\}$.
\end{enumerate}

The seal is self-verifying: strip \texttt{audit\_seal}, re-serialize, recompute
\sha, compare. Any modification to any field changes the canonical JSON and
therefore the digest.

\begin{theorem}[Seal Completeness]
The \texttt{audit\_seal.sha256} field is a collision-resistant commitment to
the complete routing decision body. Under the collision resistance of \sha,
any two distinct routing result bodies produce distinct body hashes with
overwhelming probability.
\end{theorem}

\subsection{Merkle-Chained Temporal Ledger}
\label{sec:chain}

A sealed document proves integrity of its contents but not the ordering or
completeness of a sequence of decisions. The temporal chain extends the
integrity guarantee to the entire log.

\begin{definition}[Temporal Chain]
Let $A_0 = \text{``GENESIS''}$. For the $i$-th routing decision with body
hash $B_i$, define the time anchor:
\begin{equation}
A_i = \sha\!\left(\,\texttt{encode\_utf8}(A_{i-1} \mathbin{\|} \text{``:''} \mathbin{\|} B_i)\,\right)
\label{eq:chain}
\end{equation}
where $\|$ denotes string concatenation.
\end{definition}

The \texttt{audit\_seal} block for the $i$-th decision contains:
$B_i$ (body hash), $A_i$ (\texttt{time\_anchor\_hash}),
$A_{i-1}$ (\texttt{prev\_anchor\_sha256}), and \texttt{timestamp\_utc}
(wall-clock time, per-invocation, not deterministic).

\begin{theorem}[Tamper Detection]
\label{thm:tamper}
Let $L = (e_1, \ldots, e_n)$ be a correctly-constructed chain with anchors
$(A_1, \ldots, A_n)$. If an adversary produces a modified chain $L'$ in which
entry $e_k$ is replaced by $e'_k$ with $B'_k \neq B_k$, then $A'_j \neq A_j$
for all $j \geq k$, and the chain verifier detects a mismatch at position $k$.
\end{theorem}

\begin{proof}
Since $B'_k \neq B_k$ and \sha\ is collision-resistant,
$A'_k = \sha(A_{k-1} \| \text{``:''} \| B'_k) \neq A_k$ with overwhelming
probability. By induction on $j \geq k$: each $A'_j$ depends on $A'_{j-1}
\neq A_{j-1}$ and hence $A'_j \neq A_j$. The verifier recomputes each anchor
from the stored body hashes and backlinks and detects the mismatch at position
$k$. \qed
\end{proof}

\noindent Deletion of entry $k$ produces an analogous detection: $e_{k+1}$'s
stored $A_{k-1}$ is inconsistent with the re-derived anchor from $e_{k-1}$.

\subsection{RFC~3161 External Notarization}
\label{sec:rfc3161}

The temporal chain is a single-writer structure; an adversary with full system
access could reconstruct a forged chain. RFC~3161 Trusted
Timestamping~\cite{adams2001rfc3161} closes this window.

\textbf{Notarization procedure:}
\begin{enumerate}
  \item Build a DER-encoded \texttt{TimeStampReq} (TSQ) containing: version~$=1$,
    \texttt{MessageImprint} with \sha\ OID \texttt{2.16.840.1.101.3.4.2.1} and
    the 32-byte anchor value $A_n$, and \texttt{certReq~=~TRUE}.
  \item POST the TSQ to a TSA endpoint (default: \texttt{freetsa.org}). The TSA
    returns a signed \texttt{TimeStampResponse} (TSR) containing the
    TSA-certified timestamp and certificate chain.
  \item Persist the TSR alongside the JSONL audit log.
\end{enumerate}

The TSR proves that $A_n$ existed at or before the TSA-certified time.
Forgery requires compromising the TSA, which is computationally infeasible
under standard PKI security assumptions.

\begin{corollary}[Three-Layer Compliance Proof]
\label{cor:threelayer}
The combination of body seal, temporal chain, and RFC~3161 TSR provides:
\begin{itemize}
  \item \textbf{Layer~1} (body seal): each routing decision is unmodified since sealing;
  \item \textbf{Layer~2} (temporal chain): the sequence is unmodified since construction;
  \item \textbf{Layer~3} (RFC~3161 TSR): the chain existed at or before the notarized time,
    witnessed by an independent third party.
\end{itemize}
\end{corollary}

% ─────────────────────────────────────────────────────────────────────────────
\section{Empirical Evaluation}
\label{sec:eval}

\subsection{Experimental Setup}

We evaluate TEIA on \textbf{MT-Bench}~\cite{zheng2023mtbench}: 80 multi-turn
questions across 8 categories (writing, roleplay, reasoning, math, coding,
extraction, STEM, and humanities, 10 questions each). Benchmark questions were
downloaded from the FastChat repository (Apache~2.0 license) and verified via
\sha\ digest \texttt{168710af\ldots} for reproducibility.

\textbf{Quality estimation model.}
We assign conservative quality scores:
Cloud~$= 1.00$, Hybrid~$= 0.99$, Local~$= 0.98$.
These are lower bounds on true quality for extraction tasks (where
local models match frontier models) and upper bounds for complex reasoning tasks
(where degradation is substantially greater). The reported quality retention
figure is therefore a \emph{conservative} estimate.

\textbf{Quality retention} is defined as:
\begin{equation}
Q = \frac{1}{n} \sum_{i=1}^{n} q\!\left(\text{routing}(t_i)\right)
  \;/\; 1.00
\end{equation}
where $n = 80$ and $q(\cdot)$ maps routing tier to quality score.

\textbf{Cost reduction} is defined as:
\begin{equation}
C_{\Delta} = 1 - \frac{\sum_i \text{cost}(\text{routing}(t_i))}
                      {\sum_i \text{cost}(\text{``Cloud''})}
\end{equation}
where cost is measured in estimated GPU-seconds.

\subsection{Compliance-Safe Mode Results}

\begin{table}[h]
\centering
\caption{MT-Bench results under compliance-safe mode ($\tau_{cs} = 0.20$).}
\label{tab:cs-results}
\begin{tabular}{lr}
\toprule
\textbf{Metric} & \textbf{Value} \\
\midrule
Total prompts & 80 \\
Routed to Local & 8 (10.0\%) \\
Routed to Hybrid & 0 (0.0\%) \\
Routed to Cloud & 72 (90.0\%) \\
Quality retention vs.\ 100\%-Cloud & \textbf{99.6\%} \\
Cost reduction vs.\ 100\%-Cloud & \textbf{16.3\%} \\
Routing latency, median & $< 0.1$\,ms \\
Routing latency, P99 & $< 1$\,ms \\
Throughput (single core) & $>15{,}000$ prompts/sec \\
Benchmark verdict & \textbf{PASS} (target: $\geq$95\%) \\
\bottomrule
\end{tabular}
\end{table}

The 8 prompts routed to Local are MT-Bench extraction-category questions with
$E(t) < 0.20$. All math, reasoning, coding, STEM, writing, and roleplay prompts
scored above the compliance-safe threshold and routed to Cloud.

\subsection{Max-Savings Mode Results}

\begin{table}[h]
\centering
\caption{MT-Bench results under max-savings mode ($\tau_L = 0.35$, $\tau_H = 0.65$).}
\label{tab:ms-results}
\begin{tabular}{lr}
\toprule
\textbf{Metric} & \textbf{Value} \\
\midrule
Routed to Local  & 50.0\% \\
Routed to Hybrid & 47.5\% \\
Routed to Cloud  & 2.5\% \\
Quality retention & 73.8\% \\
Cost reduction    & 98.4\% \\
\bottomrule
\end{tabular}
\end{table}

Max-savings mode achieves 98.4\% cost reduction but at 73.8\% quality retention
--- below the 95\% threshold recommended for regulated deployments.

\subsection{Performance Benchmarks}

\begin{table}[h]
\centering
\caption{Performance characteristics (i3-10100F 3.6\,GHz, single core).}
\label{tab:perf}
\begin{tabular}{lrl}
\toprule
\textbf{Metric} & \textbf{Value} & \textbf{Notes} \\
\midrule
Routing latency, median & $<0.1$\,ms & \\
Routing latency, P99 & $<1$\,ms & \\
Throughput & $>15{,}000$/sec & Single core \\
Memory footprint & $<5$\,MB & Process RSS \\
Cold-start & $<50$\,ms & Python import \\
GPU requirement & None & No model inference \\
Network calls (routing) & None & Fully offline \\
External dependencies & None & Python 3.8+ stdlib \\
\bottomrule
\end{tabular}
\end{table}

% ─────────────────────────────────────────────────────────────────────────────
\section{Regulatory Compliance Analysis}
\label{sec:regulatory}

\subsection{EU AI Act (2024)}

For High-Risk AI Systems (Annex~III), the EU AI Act imposes:

\textbf{Article~12 --- Record-keeping:} ``High-risk AI systems shall technically
allow for the automatic recording of events throughout the lifetime of the
system.'' TEIA satisfies this via the JSONL gateway audit log with SHA-256 seal
and temporal chain anchor on every routing decision.

\textbf{Article~13 --- Transparency:} TEIA's \texttt{routing\_rationale} field
provides a natural-language explanation for every decision, generated causally
at decision time.

\textbf{Annex~IV --- Technical documentation:} The fixed-version formula
(P40.0, no drift), SHA-256 seal, and temporal chain together constitute the
required technical documentation. ML-based routers cannot satisfy Annex~IV's
model change management provisions because their behavior changes at retraining.

\subsection{GDPR Article~22}

GDPR Recital~71 requires automated decisions to be ``subject to suitable
safeguards,'' including the right to explanation. TEIA's \texttt{routing\_rationale}
provides this causally, not post-hoc.

GDPR Article~5(2) accountability: the SHA-256 audit seal and temporal chain
provide the technical mechanism for demonstrating compliance.

\subsection{SEC/FINRA Algorithmic Accountability}

FINRA Rule~3110 and SEC Regulation~S-7 require algorithmic systems to maintain
records enabling audit and reconstruction. TEIA satisfies this directly: the
fixed-version formula can be re-executed from any historical prompt to produce
the identical routing result and SHA-256 seal.

\subsection{HIPAA Security Rule \S164.312(b)}

The Audit Controls standard requires mechanisms that ``record and examine
activity in information systems that contain or use ePHI.'' TEIA addresses:

\begin{itemize}
  \item \textbf{Audit log:} JSONL log with temporal chain enables post-hoc
    tamper detection.
  \item \textbf{Prompt locality:} Routing is performed in-process without
    forwarding prompt text to external services, supporting the minimum-necessary
    principle for ePHI-containing prompts.
\end{itemize}

\subsection{Basel~III / SR~11-7 Model Risk Management}

The Federal Reserve SR~11-7 guidance requires model validation including
conceptual soundness assessment and reproducible backtesting. TEIA's fixed
formula satisfies SR~11-7 directly: each feature is interpretable, the formula
is fully documented, and identical inputs produce identical outputs, enabling
straightforward backtesting.

% ─────────────────────────────────────────────────────────────────────────────
\section{Discussion and Limitations}
\label{sec:limitations}

\subsection{Heuristic Accuracy vs.\ ML Accuracy}

TEIA's formula is a structured heuristic, not a learned classifier. Its
misclassification modes are predictable:

\begin{itemize}
  \item \textbf{Over-routing to Cloud:} A short analytical-sounding prompt may
    score above $\tau_{cs}$ due to reasoning verb presence.
  \item \textbf{Under-routing to Cloud:} A long extraction prompt may score
    above $\tau_{cs}$ due to token length contribution.
\end{itemize}

Critically, these misclassifications are \emph{deterministic and auditable}:
reproducible, explainable by feature scores, and correctable by adjusting
thresholds with a protocol version increment. ML router misclassifications are
non-reproducible and cannot be corrected without retraining.

\subsection{Quality Estimation Methodology}

Our quality estimation model (Cloud~$=1.00$, Hybrid~$=0.99$, Local~$=0.98$)
is a simplification. True quality is task- and model-dependent. The reported
99.6\% quality retention should be interpreted as a conservative lower bound on
the quality fraction achievable by compliance-safe routing on the MT-Bench
distribution. A rigorous quality study would require running actual inference at
each tier.

\subsection{TSA Availability and Trust}

Production deployments should use a TSA with a documented SLA and a CA
certificate chain compatible with the relevant jurisdiction's legal framework
(e.g., EU Trust List under eIDAS). The \texttt{teia-notarize --tsa} flag
accepts any RFC~3161-compliant endpoint.

\subsection{Single-Writer Chain Architecture}

A malicious insider with write access could delete the log and reconstruct a
forged chain before an RFC~3161 checkpoint. Mitigations: frequent notarization
(after every $N$ decisions), replication to immutable write-once storage
(e.g., S3 Object Lock), or integration with a distributed ledger.

% ─────────────────────────────────────────────────────────────────────────────
\section{Conclusion}
\label{sec:conclusion}

We have presented TEIA Cognitive Router --- a deterministic, compliance-first
LLM routing system that produces cryptographically sealed, reproducible routing
decisions with sub-millisecond latency and zero ML dependencies. The system is
designed from first principles to satisfy the audit and explainability
requirements of regulated industry deployments.

The core contribution is Theorem~\ref{thm:wr}: the formal proof that the routing
formula satisfies the Write\,==\,Read invariant. This property --- which no
ML-based router can provide --- is the foundation for compliance with EU AI Act
Annex~IV, GDPR Article~22, HIPAA \S164.312(b), and SEC/FINRA record
reconstruction requirements.

The MT-Bench evaluation demonstrates that compliance-safe mode achieves 99.6\%
quality retention at 16.3\% cost reduction --- a Pareto improvement that does
not require sacrificing quality for auditability. The temporal chain extends the
integrity guarantee to entire decision sequences, and RFC~3161 external
notarization (Corollary~\ref{cor:threelayer}) closes the retroactive forgery
window.

The implementation is available under Apache~2.0 at
\texttt{pip install teia-cognitive-router}. Future work includes
domain-specific formula calibration, per-tenant chain isolation, integration
with commercial eIDAS-qualified TSAs, and a rigorous inference-level quality
study comparing actual model outputs at each routing tier.

% ─────────────────────────────────────────────────────────────────────────────
\section*{Acknowledgments}

The author thanks the lm-sys/FastChat project for releasing the MT-Bench
question set under the Apache~2.0 license. The FreeTSA.org service provides
free RFC~3161 timestamp tokens that made the notarization module practical
for open-source deployment.

% ─────────────────────────────────────────────────────────────────────────────
\bibliographystyle{plain}
\begin{thebibliography}{99}

\bibitem{adams2001rfc3161}
C.~Adams, P.~Cain, D.~Pinkas, and R.~Zuccherato.
\newblock Internet X.509 Public Key Infrastructure Time-Stamp Protocol (TSP).
\newblock RFC 3161, IETF, 2001.
\newblock \url{https://www.rfc-editor.org/rfc/rfc3161}

\bibitem{chen2023frugal}
L.~Chen, M.~Zaharia, and J.~Zou.
\newblock FrugalGPT: How to Use Large Language Models While Reducing Cost and
  Improving Performance.
\newblock \textit{arXiv:2305.05176}, 2023.

\bibitem{ding2024hybrid}
D.~Ding, A.~Mallick, C.~Wang, R.~Sim, S.~Mukherjee, O.~Ruwase, Y.~He, and
  B.~Nushi.
\newblock Hybrid LLM: Cost-Efficient and Quality-Aware Query Routing.
\newblock \textit{arXiv:2404.14618}, 2024.

\bibitem{doshi2017interpretable}
F.~Doshi-Velez and B.~Kim.
\newblock Towards a Rigorous Science of Interpretable Machine Learning.
\newblock \textit{arXiv:1702.08608}, 2017.

\bibitem{euaiact2024}
European Parliament.
\newblock Regulation (EU) 2024/1689 of the European Parliament and of the
  Council (EU AI Act).
\newblock \textit{Official Journal of the European Union}, 2024.

\bibitem{jiang2023blender}
D.~Jiang, X.~Ren, and B.~Y.~Lin.
\newblock LLM-Blender: Ensembling Large Language Models with Pairwise Ranking
  and Generative Fusion.
\newblock \textit{arXiv:2306.02561}, 2023.

\bibitem{laurie2013ct}
B.~Laurie, A.~Langley, and E.~Kasper.
\newblock Certificate Transparency.
\newblock RFC 6962, IETF, 2013.
\newblock \url{https://www.rfc-editor.org/rfc/rfc6962}

\bibitem{ong2024routellm}
I.~Ong, A.~Almahairi, V.~Wu, W.-L.~Chiang, T.~Wu, J.~E.~Gonzalez,
  M.~W.~Kadous, and I.~Stoica.
\newblock RouteLLM: Learning to Route LLMs with Preference Data.
\newblock \textit{arXiv:2406.18665}, 2024.

\bibitem{zheng2023mtbench}
L.~Zheng, W.-L.~Chiang, Y.~Sheng, S.~Zhuang, Z.~Wu, Y.~Zhuang, Z.~Lin,
  Z.~Li, D.~Li, E.~Xing, H.~Zhang, J.~E.~Gonzalez, and I.~Stoica.
\newblock Judging LLM-as-a-Judge with MT-Bench and Chatbot Arena.
\newblock \textit{arXiv:2306.05685}, 2023.

\end{thebibliography}

\end{document}
